% Options for packages loaded elsewhere
\PassOptionsToPackage{unicode}{hyperref}
\PassOptionsToPackage{hyphens}{url}
\PassOptionsToPackage{dvipsnames,svgnames,x11names}{xcolor}
%
\documentclass[
  letterpaper,
  DIV=11,
  numbers=noendperiod]{scrreprt}

\usepackage{amsmath,amssymb}
\usepackage{iftex}
\ifPDFTeX
  \usepackage[T1]{fontenc}
  \usepackage[utf8]{inputenc}
  \usepackage{textcomp} % provide euro and other symbols
\else % if luatex or xetex
  \usepackage{unicode-math}
  \defaultfontfeatures{Scale=MatchLowercase}
  \defaultfontfeatures[\rmfamily]{Ligatures=TeX,Scale=1}
\fi
\usepackage{lmodern}
\ifPDFTeX\else  
    % xetex/luatex font selection
\fi
% Use upquote if available, for straight quotes in verbatim environments
\IfFileExists{upquote.sty}{\usepackage{upquote}}{}
\IfFileExists{microtype.sty}{% use microtype if available
  \usepackage[]{microtype}
  \UseMicrotypeSet[protrusion]{basicmath} % disable protrusion for tt fonts
}{}
\makeatletter
\@ifundefined{KOMAClassName}{% if non-KOMA class
  \IfFileExists{parskip.sty}{%
    \usepackage{parskip}
  }{% else
    \setlength{\parindent}{0pt}
    \setlength{\parskip}{6pt plus 2pt minus 1pt}}
}{% if KOMA class
  \KOMAoptions{parskip=half}}
\makeatother
\usepackage{xcolor}
\setlength{\emergencystretch}{3em} % prevent overfull lines
\setcounter{secnumdepth}{5}
% Make \paragraph and \subparagraph free-standing
\makeatletter
\ifx\paragraph\undefined\else
  \let\oldparagraph\paragraph
  \renewcommand{\paragraph}{
    \@ifstar
      \xxxParagraphStar
      \xxxParagraphNoStar
  }
  \newcommand{\xxxParagraphStar}[1]{\oldparagraph*{#1}\mbox{}}
  \newcommand{\xxxParagraphNoStar}[1]{\oldparagraph{#1}\mbox{}}
\fi
\ifx\subparagraph\undefined\else
  \let\oldsubparagraph\subparagraph
  \renewcommand{\subparagraph}{
    \@ifstar
      \xxxSubParagraphStar
      \xxxSubParagraphNoStar
  }
  \newcommand{\xxxSubParagraphStar}[1]{\oldsubparagraph*{#1}\mbox{}}
  \newcommand{\xxxSubParagraphNoStar}[1]{\oldsubparagraph{#1}\mbox{}}
\fi
\makeatother


\providecommand{\tightlist}{%
  \setlength{\itemsep}{0pt}\setlength{\parskip}{0pt}}\usepackage{longtable,booktabs,array}
\usepackage{calc} % for calculating minipage widths
% Correct order of tables after \paragraph or \subparagraph
\usepackage{etoolbox}
\makeatletter
\patchcmd\longtable{\par}{\if@noskipsec\mbox{}\fi\par}{}{}
\makeatother
% Allow footnotes in longtable head/foot
\IfFileExists{footnotehyper.sty}{\usepackage{footnotehyper}}{\usepackage{footnote}}
\makesavenoteenv{longtable}
\usepackage{graphicx}
\makeatletter
\def\maxwidth{\ifdim\Gin@nat@width>\linewidth\linewidth\else\Gin@nat@width\fi}
\def\maxheight{\ifdim\Gin@nat@height>\textheight\textheight\else\Gin@nat@height\fi}
\makeatother
% Scale images if necessary, so that they will not overflow the page
% margins by default, and it is still possible to overwrite the defaults
% using explicit options in \includegraphics[width, height, ...]{}
\setkeys{Gin}{width=\maxwidth,height=\maxheight,keepaspectratio}
% Set default figure placement to htbp
\makeatletter
\def\fps@figure{htbp}
\makeatother
% definitions for citeproc citations
\NewDocumentCommand\citeproctext{}{}
\NewDocumentCommand\citeproc{mm}{%
  \begingroup\def\citeproctext{#2}\cite{#1}\endgroup}
\makeatletter
 % allow citations to break across lines
 \let\@cite@ofmt\@firstofone
 % avoid brackets around text for \cite:
 \def\@biblabel#1{}
 \def\@cite#1#2{{#1\if@tempswa , #2\fi}}
\makeatother
\newlength{\cslhangindent}
\setlength{\cslhangindent}{1.5em}
\newlength{\csllabelwidth}
\setlength{\csllabelwidth}{3em}
\newenvironment{CSLReferences}[2] % #1 hanging-indent, #2 entry-spacing
 {\begin{list}{}{%
  \setlength{\itemindent}{0pt}
  \setlength{\leftmargin}{0pt}
  \setlength{\parsep}{0pt}
  % turn on hanging indent if param 1 is 1
  \ifodd #1
   \setlength{\leftmargin}{\cslhangindent}
   \setlength{\itemindent}{-1\cslhangindent}
  \fi
  % set entry spacing
  \setlength{\itemsep}{#2\baselineskip}}}
 {\end{list}}
\usepackage{calc}
\newcommand{\CSLBlock}[1]{\hfill\break\parbox[t]{\linewidth}{\strut\ignorespaces#1\strut}}
\newcommand{\CSLLeftMargin}[1]{\parbox[t]{\csllabelwidth}{\strut#1\strut}}
\newcommand{\CSLRightInline}[1]{\parbox[t]{\linewidth - \csllabelwidth}{\strut#1\strut}}
\newcommand{\CSLIndent}[1]{\hspace{\cslhangindent}#1}

\KOMAoption{captions}{tableheading}
\makeatletter
\@ifpackageloaded{bookmark}{}{\usepackage{bookmark}}
\makeatother
\makeatletter
\@ifpackageloaded{caption}{}{\usepackage{caption}}
\AtBeginDocument{%
\ifdefined\contentsname
  \renewcommand*\contentsname{Table of contents}
\else
  \newcommand\contentsname{Table of contents}
\fi
\ifdefined\listfigurename
  \renewcommand*\listfigurename{List of Figures}
\else
  \newcommand\listfigurename{List of Figures}
\fi
\ifdefined\listtablename
  \renewcommand*\listtablename{List of Tables}
\else
  \newcommand\listtablename{List of Tables}
\fi
\ifdefined\figurename
  \renewcommand*\figurename{Figure}
\else
  \newcommand\figurename{Figure}
\fi
\ifdefined\tablename
  \renewcommand*\tablename{Table}
\else
  \newcommand\tablename{Table}
\fi
}
\@ifpackageloaded{float}{}{\usepackage{float}}
\floatstyle{ruled}
\@ifundefined{c@chapter}{\newfloat{codelisting}{h}{lop}}{\newfloat{codelisting}{h}{lop}[chapter]}
\floatname{codelisting}{Listing}
\newcommand*\listoflistings{\listof{codelisting}{List of Listings}}
\makeatother
\makeatletter
\makeatother
\makeatletter
\@ifpackageloaded{caption}{}{\usepackage{caption}}
\@ifpackageloaded{subcaption}{}{\usepackage{subcaption}}
\makeatother

\ifLuaTeX
  \usepackage{selnolig}  % disable illegal ligatures
\fi
\usepackage{bookmark}

\IfFileExists{xurl.sty}{\usepackage{xurl}}{} % add URL line breaks if available
\urlstyle{same} % disable monospaced font for URLs
\hypersetup{
  pdftitle={Desarrollo e Implementación con tecnologías libres de un Dispositivo Portátil para Evaluar el Confort Térmico, Lumínico y Acústico en Interiores},
  pdfauthor={José Ramón Hernández Aguilar},
  colorlinks=true,
  linkcolor={blue},
  filecolor={Maroon},
  citecolor={Blue},
  urlcolor={Blue},
  pdfcreator={LaTeX via pandoc}}


\title{Desarrollo e Implementación con tecnologías libres de un
Dispositivo Portátil para Evaluar el Confort Térmico, Lumínico y
Acústico en Interiores}
\author{José Ramón Hernández Aguilar}
\date{2024-12-25}

\begin{document}
\maketitle

\renewcommand*\contentsname{Table of contents}
{
\hypersetup{linkcolor=}
\setcounter{tocdepth}{2}
\tableofcontents
}

\bookmarksetup{startatroot}

\chapter*{Portada}\label{portada}
\addcontentsline{toc}{chapter}{Portada}

\markboth{Portada}{Portada}

This is a Quarto book.

To learn more about Quarto books visit
\url{https://quarto.org/docs/books}.

\bookmarksetup{startatroot}

\chapter{Introducción al Confort en Espacios y el Desarrollo del
DTHIS-C}\label{introducciuxf3n-al-confort-en-espacios-y-el-desarrollo-del-dthis-c}

\section{Importancia del Confort}\label{importancia-del-confort}

\section{Problemática de los Métodos de Medición
Tradicionales}\label{problemuxe1tica-de-los-muxe9todos-de-mediciuxf3n-tradicionales}

\section{Objetivos}\label{objetivos}

\bookmarksetup{startatroot}

\chapter{Tipos de Confort en
Espacios}\label{tipos-de-confort-en-espacios}

\section{Definición de Confort en
Espacios}\label{definiciuxf3n-de-confort-en-espacios}

\section{Confort Térmico}\label{confort-tuxe9rmico}

\section{Confort Acústico}\label{confort-acuxfastico}

\section{Confort Lumínico}\label{confort-lumuxednico}

\section{Calidad del Aire}\label{calidad-del-aire}

\bookmarksetup{startatroot}

\chapter{Diseño y Construcción del
DTHIS-C}\label{diseuxf1o-y-construcciuxf3n-del-dthis-c}

\section{Introducción al DTHIS-C}\label{introducciuxf3n-al-dthis-c}

\section{Componentes Principales del
DTHIS-C}\label{componentes-principales-del-dthis-c}

\section{Diseño del Sistema}\label{diseuxf1o-del-sistema}

\section{Medición de Variables y Metodología de
Implementación}\label{mediciuxf3n-de-variables-y-metodologuxeda-de-implementaciuxf3n}

\subsection{Temperatura del Aire}\label{temperatura-del-aire}

\subsection{Temperatura Radiante}\label{temperatura-radiante}

\subsection{Velocidad del Viento}\label{velocidad-del-viento}

\subsection{Humedad Relativa}\label{humedad-relativa}

\subsection{Concentración de CO2}\label{concentraciuxf3n-de-co2}

\subsection{Luminancia y
Deslumbramiento}\label{luminancia-y-deslumbramiento}

\subsection{Niveles de Sonido}\label{niveles-de-sonido}

\section{Calibración y Validación del
DTHIS-C}\label{calibraciuxf3n-y-validaciuxf3n-del-dthis-c}

\bookmarksetup{startatroot}

\chapter{Resultados y Análisis de las Pruebas del
DTHIS-C}\label{resultados-y-anuxe1lisis-de-las-pruebas-del-dthis-c}

\section{Presentación de los Resultados por Variables de
Confort}\label{presentaciuxf3n-de-los-resultados-por-variables-de-confort}

\section{Análisis e Interpretación de
Resultados}\label{anuxe1lisis-e-interpretaciuxf3n-de-resultados}

\bookmarksetup{startatroot}

\chapter{Discusión y Conclusiones}\label{discusiuxf3n-y-conclusiones}

\section{Resumen General de la Tesis}\label{resumen-general-de-la-tesis}

\section{Visualización de Datos}\label{visualizaciuxf3n-de-datos}

\section{Impacto del DTHIS-C en la Medición del
Confort}\label{impacto-del-dthis-c-en-la-mediciuxf3n-del-confort}

\section{Limitaciones y Propuestas de
Mejora}\label{limitaciones-y-propuestas-de-mejora}

\bookmarksetup{startatroot}

\chapter*{Bibliografía}\label{bibliografuxeda}
\addcontentsline{toc}{chapter}{Bibliografía}

\markboth{Bibliografía}{Bibliografía}

\phantomsection\label{refs}
\begin{CSLReferences}{1}{0}
\bibitem[\citeproctext]{ref-knuth84}
Knuth, Donald E. 1984. {``Literate Programming.''} \emph{Comput. J.} 27
(2): 97--111. \url{https://doi.org/10.1093/comjnl/27.2.97}.

\end{CSLReferences}

\bookmarksetup{startatroot}

\chapter*{Anexos}\label{anexos}
\addcontentsline{toc}{chapter}{Anexos}

\markboth{Anexos}{Anexos}

This is a book created from markdown and executable code.

See Knuth (1984) for additional discussion of literate programming.




\end{document}
